% Shared preamble and content for the SICVEC 2026 flyers (A4, post, story).
% Every figure here comes from the decided record (CLAUDE.md / docs/LOG.md); change it here, not per flyer.
\usepackage[T1]{fontenc}
\usepackage[utf8]{inputenc}
\usepackage[spanish]{babel}
\usepackage[default]{lato}
\usepackage{graphicx}
\graphicspath{{../../01_Propuesta/Logos/}}
\usepackage{tikz}
\usetikzlibrary{calc}
\usepackage{xcolor}
\usepackage[hidelinks]{hyperref}
\pagestyle{empty}
\setlength{\parindent}{0pt}

\definecolor{verdeosc}{RGB}{20,84,48}
\definecolor{verde}{RGB}{26,107,60}
\definecolor{verdeclaro}{RGB}{232,243,235}
\definecolor{dorado}{RGB}{212,172,13}
\definecolor{tinta}{RGB}{33,37,41}

\newcommand{\Nombre}{I Simposio Internacional de\\Ciencia Verde y Economía Circular}
\newcommand{\Sigla}{SICVEC 2026}
\newcommand{\Fechas}{19 y 20 de octubre de 2026}
\newcommand{\Lugar}{Sala de conferencias, C.C. Guacarí}
\newcommand{\Ciudad}{Sincelejo, Sucre, Colombia}
\newcommand{\Modalidad}{Modalidad híbrida}
\newcommand{\Cupos}{300 cupos en línea gratuitos para participantes externos e internacionales}
\newcommand{\Cierre}{4 de octubre de 2026, 23:59}
\newcommand{\Contacto}{insilico@unisucre.edu.co}
\newcommand{\Enlace}{Envío de resúmenes: [PENDIENTE: enlace de la plataforma]}
\newcommand{\Organiza}{Universidad de Sucre · Facultad de Educación y Ciencias · Programa de Biología · Grupo IN SILICO}
\newcommand{\Departamentos}{Departamentos: Biología y Química · Ingeniería Agroindustrial · Ingeniería Agrícola · Física}

% Institutional logos strip, scaled to a common height #1
\newcommand{\Logos}[1]{%
  \includegraphics[height=#1]{logo_unisucre.png}\hspace{1.2em}%
  \includegraphics[height=#1]{Logo_Bio.png}\hspace{1.2em}%
  \includegraphics[height=#1]{logo_agroin.png}\hspace{1.2em}%
  \includegraphics[height=#1]{logo_ingagrocola.png}\hspace{1.2em}%
  \includegraphics[height=#1]{Logo_insilico.jpeg}}
